\documentclass[12pt]{article}
\usepackage[T1]{fontenc}
\usepackage[utf8]{inputenc}
\usepackage{lmodern}
\usepackage{textcomp}
\usepackage{enumitem}
\usepackage{geometry}
\usepackage{graphicx}
\usepackage{amsmath, amssymb}
\geometry{margin=1in}
\begin{document}
\begin{center}
Physics - Graduate Level Assessment 14 \
Total Marks: 40 \
Answer all questions.
\end{center}

\section*{Multiple Choice (10 marks)}
\begin{enumerate}[label=\arabic*.]
\item Calculate $\Delta G$ for the isothermal expansion of $2.25 \mathrm{~mol}$ of an ideal gas at $325 \mathrm{~K}$ from an initial pressure of 12.0 bar to a final pressure of 2.5 bar.
\begin{enumerate}[label=(\alph*)]
    \item -7.89 $10^3 \mathrm{~J}$
    \item -12.75 $10^3 \mathrm{~J}$
    \item -8.01 $10^3 \mathrm{~J}$
    \item -10.02 $10^3 \mathrm{~J}$
    \item -11.32 $10^3 \mathrm{~J}$
\end{enumerate}
\item Why are there no impact craters on the surface of Io?
\begin{enumerate}[label=(\alph*)]
    \item Any craters that existed have been eroded through the strong winds on Io's surface.
    \item Io is protected by a layer of gas that prevents any objects from reaching its surface.
    \item The gravitational pull of nearby moons prevent objects from hitting Io.
    \item Io's magnetic field deflects incoming asteroids, preventing impact craters.
    \item Io did have impact craters but they have all been buried in lava flows.
\end{enumerate}
\item An elevator is rising. In order to determine the tension in the cables pulling it, which values (for the elevator) would you need to know?
\begin{enumerate}[label=(\alph*)]
    \item Mass, velocity, and acceleration
    \item Mass, velocity, height, and acceleration
    \item Mass and velocity
    \item Velocity, height, and acceleration
    \item Velocity and acceleration
\end{enumerate}
\item Say the pupil of your eye has a diameter of 5 mm and you have a telescope with an aperture of 50 cm. How much more light can the telescope gather than your eye?
\begin{enumerate}[label=(\alph*)]
    \item 1000 times more
    \item 50 times more
    \item 5000 times more
    \item 500 times more
    \item 10000 times more
\end{enumerate}
\item When an element ejects a beta particle, the atomic number of that element
\begin{enumerate}[label=(\alph*)]
    \item reduces by 1
    \item increases by 3
    \item increases by 1
    \item reduces by half
    \item increases by 4
\end{enumerate}
\end{enumerate}

\section*{True / False (10 marks)}
\textit{Answer True or False}
\begin{enumerate}[label=\arabic*.]
\setcounter{enumi}{5}
\item True or False: Terrestrial radiation has lower polarization and frequency compared to radiation from the Sun.
\item True or False: The proximity of a planet to other celestial bodies is the primary factor determining its volcanic and tectonic history.
\item True or False: The Fermi energy for copper is typically around 4.56 eV.
\item True or False: The power of the lens is calculated to be 3 diopters based on the given object and image distances.
\item True or False: Pluto has one medium-sized satellite and two small satellites orbiting it.
\end{enumerate}

\section*{Long Form Response (20 marks)}
\textit{Answer each question with a clear and complete explanation.}
\begin{enumerate}[label=\arabic*.]
\setcounter{enumi}{10}
\item Calculate the magnitude of the gravitational force acting on a spherical water drop with a diameter of 1.20 μm suspended in a downward-directed electric field of 462 N/C. Discuss the implications of your findings.
\item A free particle with a mass of 20 grams is initially at rest. Analyze the effects of applying a 100 dyne force for 10 seconds on its kinetic energy, detailing your calculations and underlying principles.
\end{enumerate}

\vfill
\noindent\textit{End of Paper}
\end{document}